\subsection{Numerical check of the theory}\label{sec:numcheck}
The design effects of this section are first-order (linearised) quantities, so we checked them against Monte Carlo experiments on the eleven real pools and on synthetic pools in which the disagreement rate is controlled (Supplementary Section~\SIref{sec:simulation}). Three findings matter for the rest of the paper. First, the linearised design effect predicts the efficiency of the implemented design to within Monte Carlo error: over all target--budget pairs on the real pools the ratio of observed to predicted efficiency has median 1.00 (10\%--90\% range 0.83--1.22). Second, the closed-form expression of Corollary~\SIref{prop:deff} is informative but too optimistic for nine models. It is accurate when the strata separate the agreements of a pair from its disagreements (two models), and captures the $1/\rho$ scaling and the effect of the size of the contrast; with nine models it overstates the gain on real pools by a median factor of about two, because the count $K_i$ locates any disagreement but not that of a particular pair, and the bounds of Supplementary Section~\SIref{si:theory-extra} describe the gap. We therefore report the closed form as an illustration of the mechanism rather than as a forecast. Third, for the difference estimator the linearised residual variance predicts its gain over the plain estimator within Monte Carlo error (median ratio of observed to predicted gain 0.98 over 34 pool--budget cells, Supplementary Section~\SIref{sec:sim-dif}); the gain is moderate (median 1.47) rather than a doubling, and it is largest in the law-school setting, where the label-free forecast underestimates it. An uninformative or inverted proxy leaves the allocation almost unchanged, and a proxy of the wrong sign makes the difference estimator worse than the plain one for the error contrast by a factor of about 1.3, as the theory predicts.
